\documentclass[10pt]{article}

\usepackage[margin=1in]{geometry}
\usepackage[numbers,sort&compress]{natbib}
\usepackage{booktabs}
\usepackage{amsmath}
\usepackage{amsfonts}
\usepackage{hyperref}

\newcommand{\method}{HFRVLA}
\newcommand{\smolvla}{SmolVLA}
\newcommand{\abase}{a_{\mathrm{base}}}
\newcommand{\afinal}{a_{\mathrm{final}}}
\newcommand{\deltah}{\hat{\Delta a}}

\title{\method: Wrist-Centric Fast Residual Correction for Frozen VLA Action Chunks}
\author{Anonymous}
\date{}

\begin{document}
\maketitle

\begin{abstract}
Vision-language-action policies often amortize inference by predicting action chunks, but executing a chunk open-loop can make the action stream stale when the robot or scene deviates from the observation that produced the chunk. We study whether a small wrist-camera residual module can improve chunk execution without fine-tuning the base VLA. \method{} wraps a frozen LIBERO-adapted \smolvla{} slow planner and trains only a fast residual head that consumes wrist DINOv3 patches, robot state, the frozen base action, chunk position, and slow-planner context. At deployment, the robot executes $\afinal = \abase + \alpha \cdot \mathrm{clip}(\deltah)$.
In current LIBERO simulation evidence, \method{} improves matched action-step settings over the frozen \smolvla{} baseline at short execution horizons, but long matched chunks remain difficult and require residual calibration. These results position \method{} as a lightweight, LeRobot-native testbed for execution-time correction of frozen VLA chunks rather than as a completed solution to long-horizon open-loop drift.
\end{abstract}

\section{Introduction}

Large vision-language-action (VLA) policies can use language and visual context to produce robot actions, but high-capacity policies are expensive to run at every control tick. A common deployment compromise is action chunking: the policy predicts a sequence of future actions, and the robot executes that sequence for several steps before replanning. Chunking improves inference amortization and temporal consistency, but it also weakens closed-loop reactivity. When the scene, object pose, contact condition, or robot state changes during chunk execution, later actions can be based on an observation that is no longer current.

Recent work makes this timing problem explicit. A2C2 introduces a correction head that runs every control step and adds time-aware residual corrections to frozen VLA action chunks using the latest observation and base-policy context~\citep{sendai2025a2c2}. DynamicVLA addresses the same perception--execution gap from a scheduler and model-design perspective, using continuous inference and latent-aware action streaming to reduce temporal misalignment~\citep{xie2026dynamicvla}. These directions suggest that action chunk deployment should be evaluated as a closed-loop timing problem, not only as a supervised prediction problem.

\method{} asks a narrower question: how far can a wrist-centric fast residual module go when the slow VLA is kept frozen? We use `HuggingFaceVLA/smolvla\_libero' as the frozen slow planner~\citep{shukor2025smolvla}, precompute the slow planner outputs and DINOv3 wrist features~\citep{simeoni2025dinov3}, and train a lightweight correction head through the LeRobot workflow~\citep{cadene2026lerobot}. The per-step visual correction signal is wrist-camera based; the policy is not wrist-only, because the fast module also receives $\abase$, robot state, chunk position, and slow-planner latent context.

The current paper makes four claims:
\begin{enumerate}
  \item \textbf{Problem framing.} Frozen VLA action chunks should be evaluated under matched planning and execution horizons, because default every-step replanning is not a matched baseline for chunk execution.
  \item \textbf{Method.} A small wrist-camera residual module can be trained offline from cached slow-planner features and deployed as $\abase + \alpha \cdot \mathrm{clip}(\deltah)$ without updating the slow VLA.
  \item \textbf{Artifact.} The implementation is LeRobot-native: it includes dataset recording, fast-cache construction, training, checkpoint packaging, and an evaluation registry.
  \item \textbf{Evidence and limitation.} Current results show useful gains at short execution horizons, but the long-horizon setting remains fragile and motivates delayed-chunk targets, latency features, and deployment-aligned residual losses.
\end{enumerate}

\section{Method}

\subsection{Frozen Slow Planner}

The slow planner is a frozen LIBERO-adapted \smolvla{} checkpoint. Given top and wrist observations, state, and task text, the planner outputs an action chunk. During dataset recording we store the per-step base action $\abase \in \mathbb{R}^7$, a normalized chunk index, and pooled slow-planner context vectors. The current local feature contract stores $z_{\mathrm{goal}} \in \mathbb{R}^{960}$ and $z_{\mathrm{phase}} \in \mathbb{R}^{480}$.

The slow planner is not constructed during the main fast-module training loop. Instead, HFRVLA records a custom LeRobotDataset v3 and optionally converts it into a compact frame-level fast-cache. This design keeps the provenance of $\abase$, $z_{\mathrm{goal}}$, and $z_{\mathrm{phase}}$ fixed. If the slow planner changes, the HFRVLA dataset must be re-recorded.

\subsection{Fast Wrist Residual Module}

The trainable fast path receives the current wrist DINOv3 patch features, robot state, $\abase$, chunk index, slow-planner context, and the previous correction. It predicts a full 7D residual $\deltah$. The deployment merge is
\begin{equation}
  \afinal = \abase + \alpha \cdot \mathrm{clip}(\deltah, -\delta_{\max}, \delta_{\max}).
\end{equation}
The current mainline intentionally excludes the older learned gate, contact auxiliary head, and GRU-centered design. This keeps the first paper claim interpretable: any improvement comes from the residual correction and its conditioning, not from a separate intervention policy.

\subsection{Training Objective}

The implemented baseline trains a residual target against the expert action:
\begin{equation}
  \Delta a^* = a_{\mathrm{expert}} - \abase.
\end{equation}
The current code supports final-action alignment terms that compare $\abase + \alpha \cdot \mathrm{clip}(\deltah)$ directly to $a_{\mathrm{expert}}$, along with residual and clip monitoring. The ongoing diagnosis suggests that the next important target is delayed chunk-age supervision: the fast module should learn to correct the $k$-th action from a chunk generated at $t-k$, not only the simultaneous base action at $t$.

\section{Experimental Setup and Evidence}

We evaluate on LIBERO~\citep{liu2023libero} through `lerobot-eval'. The dataset used for fast training is a recorded HFRVLA LeRobotDataset v3 with 1,693 episodes and 273,465 frames. The policy checkpoint used for current FWR-v2 sweeps is packaged as `checkpoints/hfrvla\_fwr\_chunk\_generated\_seq2\_b1024p3\_50k\_packaged'. The evaluation registry is the source of truth for completed paper-facing comparisons.

\begin{table}[t]
  \centering
  \caption{Registry-backed evidence snapshot, seed 42, LIBERO Spatial + Object combined, 100 episodes per policy/setting unless noted. Planning denotes the slow planner chunk length; execution/replan denotes how many actions are executed before replanning.}
  \begin{tabular}{llccc}
    \toprule
    Sweep & Policy & Planning & Execution/replan & Success \\
    \midrule
    action-step & \method{} 30k, $\alpha=0.5$ & 50 & 2 & 85/100 \\
    action-step & \smolvla{} & 50 & 2 & 82/100 \\
    action-step & \method{} 30k, $\alpha=0.5$ & 50 & 8 & 81/100 \\
    action-step & \smolvla{} & 50 & 8 & 77/100 \\
    matched chunk & \method{} 30k, $\alpha=0.5$ & 8 & 8 & 83/100 \\
    matched chunk & \smolvla{} & 4 & 4 & 79/100 \\
    matched chunk & \method{} 30k, $\alpha=0.5$ & 50 & 50 & 46/100 \\
    matched chunk & \smolvla{} & 50 & 50 & 43/100 \\
    \bottomrule
  \end{tabular}
  \label{tab:evidence}
\end{table}

Table~\ref{tab:evidence} supports the conservative claim that wrist residual correction helps in some matched chunk-execution regimes, especially when execution/replan is short. It also shows the main unresolved problem: when planning and execution both remain at 50 steps, both \method{} and \smolvla{} degrade sharply.

\paragraph{Active calibration sweep.}
At the time of this draft, an expanded $n\_action\_steps=50$ LIBERO-Spatial alpha/clip sweep is running with 50 episodes per setting and eval batch size 3. Completed early rows for $\alpha=0.25$ are $23/50$ at $\delta_{\max}=0.05$, $23/50$ at $0.10$, $24/50$ at $0.15$, $24/50$ at $0.18$, and $27/50$ at $0.20$. The $\delta_{\max}=0.22$ row is currently running. These rows are trend diagnostics, not final paper claims.

\section{Related Work}

\paragraph{Action chunk correction.}
A2C2 is the closest conceptual anchor. It corrects frozen VLA chunks at every control step using the latest observation, base action, chunk-position feature, and base-policy context~\citep{sendai2025a2c2}. \method{} follows the same residual-on-frozen-chunk philosophy but makes a narrower design choice: high-frequency visual correction comes from the wrist camera path, while global task context is inherited from the frozen slow planner.

\paragraph{Latency-aware VLA execution.}
DynamicVLA frames dynamic manipulation as a perception--execution gap and addresses it with continuous inference and action streaming~\citep{xie2026dynamicvla}. \method{} does not redesign the VLA or scheduler as its main contribution; instead, it studies whether a small residual head can compensate for some chunk staleness while leaving \smolvla{} frozen.

\paragraph{Slow-fast robot policies.}
Reactive Diffusion Policy uses a slow-fast policy structure for contact-rich manipulation with visual-tactile feedback~\citep{xue2025rdp}. \method{} shares the slow-fast motivation but uses only wrist vision as the fast visual signal, avoiding additional tactile sensors in the current system.

\paragraph{Robot learning infrastructure.}
LeRobot provides the data, training, and evaluation stack used by this project~\citep{cadene2026lerobot}. LIBERO provides the manipulation benchmark and task suites~\citep{liu2023libero}. DINOv3 supplies the frozen dense wrist visual features used by the residual module~\citep{simeoni2025dinov3}.

\section{Limitations and Next Steps}

The current evidence is simulation-only and should not be framed as a real-hardware validation. Long matched chunks remain the main failure case. Current training also uses a simultaneous residual target, while deployment often requires correcting stale chunk actions generated from older observations. The next method iteration should therefore prioritize delayed chunk-age residual targets, explicit age/latency features, and deployment-aligned losses that include the same $\alpha$ and clipping used at inference.

The wrist-centric design is also under-conditioned for tasks where the wrist image does not contain the target or global spatial relation. The correct claim is not that the whole policy is wrist-only; the claim is that per-step new visual evidence is wrist-based while the slow planner supplies task context.

\section{Conclusion}

\method{} is a LeRobot-native testbed for fast wrist residual correction on frozen VLA action chunks. Current results justify the problem and show modest gains in matched short-execution settings, while also exposing calibration and long-horizon weaknesses. This makes the paper direction strongest as a careful study of frozen VLA chunk correction, not as an over-claimed universal improvement over VLA replanning.

\bibliographystyle{plainnat}
\bibliography{references}

\end{document}
